\chapter{2007年四川卷（文）}
\section{单选题}
\begin{problem}
设集合 \(M = \{4, 5, 6, 8\}\)，集合 \(N = \{3, 5, 7, 8\}\)，那么 \(M \cup N =\)（\quad）

\choices
    {\( \{3, 4, 5, 6, 7, 8\} \)}
    {\( \{5, 8\} \)}
    {\( \{3, 5, 7, 8\} \)}
    {\( \{4, 5, 6, 8\} \)}
\end{problem}

\begin{answer}
A
\end{answer}

\begin{solution}
集合并集包含 \(M\) 与 \(N\) 中出现过的所有元素，去掉重复元素 \(5\)，\(8\)，得
\[
M\cup N=\{3,4,5,6,7,8\}
\]
所以选 A.
\end{solution}

\begin{problem}
函数 \( f(x) = 1 + \log_2 x \) 与 \( g(x) = 2^{-x + 1} \) 在同一直角坐标系下的图象大致是（\quad）

\choices
    {\choicebitmap[width=0.15\paperwidth]{img/2007/sichuan_liberal/q2_optA.png}}
    {\choicebitmap[width=0.15\paperwidth]{img/2007/sichuan_liberal/q2_optB.png}}
    {\choicebitmap[width=0.15\paperwidth]{img/2007/sichuan_liberal/q2_optC.png}}
    {\choicebitmap[width=0.15\paperwidth]{img/2007/sichuan_liberal/q2_optD.png}}
\end{problem}

\begin{answer}
C
\end{answer}

\begin{solution}
函数 \(f(x)=1+\log_{2}x\) 的定义域为 \(x>0\)，且 \(f'(x)=\dfrac{1}{x\ln2}>0\)，所以其图象在定义域内单调递增，并经过点 \((1,1)\)，且 \(f\left(\dfrac12\right)=0\)，\(f(2)=2\). 函数 \(g(x)=2^{1-x}\) 满足 \(g'(x)=-(\ln2)2^{1-x}<0\)，所以其图象单调递减，也经过点 \((1,1)\)，且 \(g(0)=2\)，\(g(2)=\dfrac12\). 因此两图象在 \((1,1)\) 相交，且一个递增、一个递减，结合 \(f\) 的横截距为 \(\dfrac12\)，符合题意的图象为 C.
\end{solution}

\begin{problem}
某商场买来一车苹果，从中随机抽取了 \(10\) 个苹果，其重量（单位：克）分别为：\(150\)，\(152\)，\(153\)，\(149\)，\(148\)，\(146\)，\(151\)，\(150\)，\(152\)，\(147\)，由此估计这车苹果单个重量的期望值是（\quad）

\choices
    {\(150.2\text{ 克}\)}
    {\(149.8\text{ 克}\)}
    {\(149.4\text{ 克}\)}
    {\(147.8\text{ 克}\)}
\end{problem}

\begin{answer}
B
\end{answer}

\begin{solution}
用这 \(10\) 个苹果重量的平均数估计单个苹果重量的期望值，有
\[
\frac{150+152+153+149+148+146+151+150+152+147}{10}
 =\frac{1498}{10}=149.8\text{ 克}
\]
所以选 B.
\end{solution}

\begin{problem}
如图， \(ABCD - A_{1}B_{1}C_{1}D_{1}\) 为正方体，下面结论错误的是（\quad）

\bitmapfigure[float=inline,max width=4.6cm]{img/2007/sichuan_liberal/q4_fig1.png}
\FloatBarrier

\choices
    {\(BD \parallel\) 平面 \(CB_{1}D_{1}\)}
    {\(AC_{1} \perp BD\)}
    {\(AC_{1} \perp\) 平面 \(CB_{1}D_{1}\)}
    {异面直线 \(AD\) 与 \(CB_{1}\) 角为 \(60^{\circ}\)}
\end{problem}

\begin{answer}
D
\end{answer}

\begin{solution}
建立空间直角坐标系，令 \(A=(0,0,0)\)，\(B=(1,0,0)\)，\(C=(1,1,0)\)，\(D=(0,1,0)\)，以及
\(A_{1}=(0,0,1)\)，\(B_{1}=(1,0,1)\)，\(C_{1}=(1,1,1)\)，\(D_{1}=(0,1,1)\).
平面 \(CB_{1}D_{1}\) 的方程为 \(x+y+z=2\)，而直线 \(BD\) 上各点满足 \(x+y+z=1\)，故 \(BD\parallel\) 平面 \(CB_{1}D_{1}\). 又 \(\overrightarrow{AC_{1}}=(1,1,1)\)，\(\overrightarrow{BD}=(-1,1,0)\)，
所以 \(\overrightarrow{AC_{1}}\cdot\overrightarrow{BD}=0\)，从而 \(AC_{1}\perp BD\). 平面 \(CB_{1}D_{1}\) 的法向量为 \((1,1,1)\)，故 \(AC_{1}\perp\) 平面 \(CB_{1}D_{1}\). 最后，取 \(\overrightarrow{AD}=(0,1,0)\)，\(\overrightarrow{CB_{1}}=(0,-1,1)\)，则
\[
\cos\angle(AD,CB_{1})=\frac{|(0,1,0)\cdot(0,-1,1)|}{1\cdot\sqrt2}=\frac{1}{\sqrt2}
\]
故两异面直线所成的角为 \(45^{\circ}\)，不是 \(60^{\circ}\). 所以错误的结论是 D.
\end{solution}

\begin{problem}
如果双曲线 \(\frac{x^2}{4} - \frac{y^2}{2} = 1\) 上一点 \(P\) 到双曲线右焦点的距离是 \(2\)，那么点 \(P\) 到 \(y\) 轴的距离是（\quad）

\choices
    {\(\frac{4\sqrt{6}}{3}\)}
    {\(\frac{2\sqrt{6}}{3}\)}
    {\(2\sqrt{6}\)}
    {\(2\sqrt{3}\)}
\end{problem}

\begin{answer}
A
\end{answer}

\begin{solution}
双曲线的焦半距为 \(c=\sqrt{4+2}=\sqrt6\)，右焦点为 \(F=(\sqrt6,0)\). 设 \(P=(x,y)\)，由双曲线方程得
\[
y^{2}=\frac{x^{2}}{2}-2
\]
又因为 \(PF=2\)，所以
\[
(x-\sqrt6)^{2}+y^{2}=4
\]
代入上式，得
\[
x^{2}-2\sqrt6x+6+\frac{x^{2}}2-2=4
\]
从而 \(x\left(\frac32x-2\sqrt6\right)=0\).
若 \(x=0\)，则 \(y^{2}=-2\)，不可能；因此 \(x=\dfrac{4\sqrt6}{3}>0\). 点 \(P\) 到 \(y\) 轴的距离为 \(|x|=\dfrac{4\sqrt6}{3}\)，所以选 A.
\end{solution}

\begin{problem}
设球 \(O\) 的半径是 \(1\)，\(A\)，\(B\)，\(C\) 是球面上三点. 已知 \(A\) 到 \(B\)，\(C\) 两点的球面距离都是 \(\frac{\pi}{2}\)，且二面角 \(B - OA - C\) 的大小为 \(\frac{\pi}{3}\)，则从 \(A\) 点沿球面经 \(B\)，\(C\) 两点再回到 \(A\) 点的最短距离是（\quad）

\bitmapfigure[max width=6.0cm]{img/2007/sichuan_liberal/q6_fig1.png}

\choices
    {\(\frac{7\pi}{6}\)}
    {\(\frac{5\pi}{4}\)}
    {\(\frac{4\pi}{3}\)}
    {\(\frac{3\pi}{2}\)}
\end{problem}

\begin{answer}
C
\end{answer}

\begin{solution}
球的半径为 \(1\)，且 \(A\) 到 \(B\)，\(C\) 的球面距离均为 \(\dfrac\pi2\)，所以
\[
\angle AOB=\angle AOC=\frac\pi2
\]
由于 \(OB\)，\(OC\) 都垂直于 \(OA\)，二面角 \(B-OA-C\) 的平面角就是 \(\angle BOC\)，故 \(\angle BOC=\dfrac\pi3\). 半径为 \(1\) 时，球面最短弧长等于相应的圆心角，因此所求最短距离为
\[
\frac\pi2+\frac\pi3+\frac\pi2=\frac{4\pi}{3}
\]
任意一条经 \(A\)，\(B\)，\(C\)，\(A\) 的球面路径，其三段长度分别不小于这三个点对之间的球面距离；取三段对应的大圆短弧时可以同时达到这些下界，所以最短距离为 \(\frac{4\pi}{3}\)，选 C.
\end{solution}

\begin{problem}
等差数列 \(\{a_{n}\}\) 中，\(a_{1} = 1\)，\(a_{3} + a_{5} = 14\)，其前 \(n\) 项和 \(S_{n} = 100\)，则 \(n =\)（\quad）

\choices
    {\(9\)}
    {\(10\)}
    {\(11\)}
    {\(12\)}
\end{problem}

\begin{answer}
B
\end{answer}

\begin{solution}
设等差数列的公差为 \(d\)，则
\[
a_{3}+a_{5}=(1+2d)+(1+4d)=14
\]
所以 \(d=2\). 因而
\[
S_{n}=\frac{n}{2}\bigl(2a_{1}+(n-1)d\bigr)
 =\frac{n}{2}\bigl(2+2(n-1)\bigr)=n^{2}
\]
由 \(S_{n}=100\) 且 \(n\) 为正整数，得 \(n^{2}=100\)，所以 \(n=10\)，选 B.
\end{solution}

\begin{problem}
设 \(A(a,1)\)，\(B(2,b)\)，\(C(4,5)\) 为坐标平面上三点，\(O\) 为坐标原点. 若 \(\overrightarrow{OA}\) 与 \(\overrightarrow{OB}\) 在 \(\overrightarrow{OC}\) 上的投影相同，则 \(a\) 与 \(b\) 满足的关系式为（\quad）

\choices
    {\(4a - 5b = 3\)}
    {\(5a - 4b = 3\)}
    {\(4a + 5b = 14\)}
    {\(5a + 4b = 14\)}
\end{problem}

\begin{answer}
A
\end{answer}

\begin{solution}
两个向量在同一非零向量 \(\overrightarrow{OC}\) 上的投影相同，等价于它们与 \(\overrightarrow{OC}\) 的内积相同. 因为 \(\overrightarrow{OA}=(a,1)\)，\(\overrightarrow{OB}=(2,b)\)，\(\overrightarrow{OC}=(4,5)\)，
所以
\[
4a+5=\overrightarrow{OA}\cdot\overrightarrow{OC}
 =\overrightarrow{OB}\cdot\overrightarrow{OC}=8+5b
\]
整理得 \(4a-5b=3\)，所以选 A.
\end{solution}

\begin{problem}
用数字 \(1\)，\(2\)，\(3\)，\(4\)，\(5\) 可以组成没有重复数字，并且比 \(20000\) 大的五位偶数共有（\quad）

\choices
    {\(48\) 个}
    {\(36\) 个}
    {\(24\) 个}
    {\(18\) 个}
\end{problem}

\begin{answer}
B
\end{answer}

\begin{solution}
五位偶数的个位只能是 \(2\) 或 \(4\). 当个位为 \(2\) 时，首位为 \(3\)，\(4\) 或 \(5\)，有 \(3\) 种选法；首位确定后，中间三位有 \(3!\) 种排法，共 \(3\times3!=18\) 个. 当个位为 \(4\) 时，首位为 \(2\)，\(3\) 或 \(5\)，同样有 \(18\) 个. 因而总数为
\[
18+18=36
\]
所以选 B.
\end{solution}

\begin{problem}
已知抛物线 \( y = -x^2 + 3 \) 上存在关于直线 \( x + y = 0 \) 对称的相异两点 \(A\)，\(B\)，则 \( |AB| \) 等于（\quad）

\choices
    {\(3\)}
    {\(4\)}
    {\(3 \sqrt{2}\)}
    {\(4\sqrt{2}\)}
\end{problem}

\begin{answer}
C
\end{answer}

\begin{solution}
设抛物线上的一点为 \(A=(t,3-t^{2})\). 关于直线 \(x+y=0\) 对称时，点 \(A\) 的对称点为
\[
B=(t^{2}-3,-t)
\]
要使 \(B\) 仍在抛物线上，需满足
\[
-t=-\left(t^{2}-3\right)^{2}+3
\]
从而
\[
t^{4}-6t^{2}-t+6=(t-1)(t+2)(t^{2}-t-3)=0
\]
若 \(t^{2}-t-3=0\)，则 \(t^{2}=t+3\)，从而 \(A=(t,-t)=B\)，不符合两点相异. 因此 \(t=1\) 或 \(t=-2\)，对应的两点为 \(A=(1,2)\)，\(B=(-2,-1)\).
所以
\[
|AB|=\sqrt{(1+2)^{2}+(2+1)^{2}}=3\sqrt2
\]
所以选 C.
\end{solution}

\begin{problem}
某公司有 \(60\text{ 万元}\) 资金，计划投资甲，乙两个项目，按要求对项目甲的投资不小于对项目乙投资的 \(\frac{2}{3}\) 倍，且对每个项目的投资不能低于 \(5\text{ 万元}\)，对项目甲每投资 \(1\text{ 万元}\) 可获得 \(0.4\text{ 万元}\) 的利润，对项目乙每投资 \(1\text{ 万元}\) 可获得 \(0.6\text{ 万元}\) 的利润，该公司正确规划投资后，在这两个项目上共可获得的最大利润为（\quad）

\choices
    {\(36\text{ 万元}\)}
    {\(31.2\text{ 万元}\)}
    {\(30.4\text{ 万元}\)}
    {\(24\text{ 万元}\)}
\end{problem}

\begin{answer}
B
\end{answer}

\begin{solution}
设投资甲、乙项目的金额分别为 \(x\)，\(y\) 万元，则
\(x+y\leqslant60\)，\(x\geqslant\frac23y\)，\(x\geqslant5\)，\(y\geqslant5\)
由于增加甲项目的投资会增加利润且不破坏其他约束，所以最优方案必有 \(x+y=60\). 于是 \(x=60-y\)，由 \(x\geqslant\dfrac23y\) 得
\(60-y\geqslant\frac23y\)，\(y\leqslant36\)
总利润为
\[
0.4x+0.6y=0.4(60-y)+0.6y=24+0.2y
\]
它随 \(y\) 增大而增大，因此在 \(y=36\)，\(x=24\) 时取得最大值
\[
0.4\times24+0.6\times36=31.2\text{ 万元}
\]
所以选 B.
\end{solution}

\begin{problem}
如图，\(l_{1}\)，\(l_{2}\)，\(l_{3}\) 是同一平面内的三条平行直线，\(l_{1}\) 与 \(l_{2}\) 间的距离是 \(1\)，\(l_{2}\) 与 \(l_{3}\) 间的距离是 \(2\)，正三角形 \(ABC\) 的三顶点分别在 \(l_{1}\)，\(l_{2}\)，\(l_{3}\) 上，则 \(\triangle ABC\) 的边长是（\quad）

\bitmapfigure[max width=0.54\linewidth]{img/2007/sichuan_liberal/q12_fig1.png}

\choices
    {\(2\sqrt{3}\)}
    {\(\frac{4\sqrt{6}}{3}\)}
    {\(\frac{3 - \sqrt{7}}{4}\)}
    {\(\frac{2 - \sqrt{21}}{3}\)}
\end{problem}

\begin{answer}
\(\frac{2\sqrt{21}}{3}\)（按题面条件计算，原卷选项有误）.
\end{answer}

\begin{solution}
取三条平行线分别为 \(y=1\)，\(y=0\)，\(y=-2\)，设 \(A=(x_{A},1)\)，\(B=(x_{B},0)\)，\(C=(x_{C},-2)\)，
并记 \(u=x_{A}-x_{B}\)，\(v=x_{B}-x_{C}\). 设正三角形边长为 \(s\)，则
\[
s^{2}=u^{2}+1=v^{2}+4=(u+v)^{2}+9
\]
由 \(AB=BC\) 和 \(AB=AC\) 分别得
\[
\begin{cases}
u^{2}-v^{2}=3,\\
u^{2}+1=(u+v)^{2}+9
\end{cases}
\]
从而
\[
\begin{cases}
u^{2}-v^{2}=3,\\
2uv+v^{2}+8=0
\end{cases}
\]
由 \(u^{2}=v^{2}+3\)，对第二个等式平方，得
\(4u^{2}v^{2}=(v^{2}+8)^{2}\)，\(4v^{2}(v^{2}+3)=(v^{2}+8)^{2}\)
整理得
\[
3v^{4}-4v^{2}-64=0
\]
令 \(z=v^{2}\geqslant0\)，则 \(3z^{2}-4z-64=0\)，故
\[
z=\frac{4\pm\sqrt{16+768}}{6}=\frac{4\pm28}{6}
\]
舍去负根，得 \(v^{2}=\dfrac{16}{3}\)，从而 \(u^{2}=\dfrac{25}{3}\). 此时可取 \(u=\dfrac5{\sqrt3}\)，\(v=-\dfrac4{\sqrt3}\)，满足 \(2uv+v^{2}+8=0\)，故平方没有引入不合条件的结果. 因而
\[
s^{2}=v^{2}+4=\frac{16}{3}+4=\frac{28}{3}
\]
所以 \(s=\sqrt{\dfrac{28}{3}}=\dfrac{2\sqrt{21}}{3}\). 该值与原卷第四项印刷的 \(\dfrac{2-\sqrt{21}}{3}\) 不一致，故按题面条件计算，原卷选项与题面条件不相容.
\end{solution}

\section{填空题}
\begin{problem}
\(\left(x - \frac{1}{x}\right)^n\) 的展开式中的第 \(5\) 项为常数项，那么正整数 \(n\) 的值是\fillinblank{}.
\end{problem}

\begin{answer}
\(8\)
\end{answer}

\begin{solution}
二项式展开式的第 \(5\) 项对应于通项中的 \(k=4\)，其通项为
\[
T_{k+1}=\mathrm C_n^k x^{n-k}\left(-\frac{1}{x}\right)^k=(-1)^k\mathrm C_n^k x^{n-2k}
\]
因此第 \(5\) 项为 \((-1)^4\mathrm C_n^4x^{n-8}\). 它为常数项，当且仅当
\[
n-8=0
\]
所以 \(n=8\).
\end{solution}

\begin{problem}
如图，在正三棱柱 \(ABC - A_{1}B_{1}C_{1}\) 中，侧棱长为 \(\sqrt{2}\)，底面三角形的边长为 1，则 \(BC_{1}\) 与侧面 \(ACC_{1}A_{1}\) 所成的角是\fillinblank{}.
\FigureLayoutDeclare{img/2007/sichuan_liberal/q14_fig1.png}{5.0cm}{5.999bp 5.999bp 5.999bp 5.999bp}
\bitmapfigure[max width=6.0cm]{img/2007/sichuan_liberal/q14_fig1.png}
\end{problem}

\begin{answer}
\(\dfrac{\pi}{6}\)
\end{answer}

\begin{solution}
建立空间直角坐标系，令 \(A=(0,0,0)\)，\(C=(1,0,0)\)，\(B=\left(\dfrac12,\dfrac{\sqrt3}{2},0\right)\)，\(C_{1}=(1,0,\sqrt2)\). 则侧面 \(ACC_{1}A_{1}\) 的方程为 \(y=0\)，其法向量可取 \(\bs{n}=(0,1,0)\)，且 \(\overrightarrow{BC_{1}}=\left(\frac12,-\frac{\sqrt3}{2},\sqrt2\right)\)，
\[
\left|\overrightarrow{BC_{1}}\right|=\sqrt{\frac14+\frac34+2}=\sqrt3
\]
设直线 \(BC_{1}\) 与平面 \(ACC_{1}A_{1}\) 所成的角为 \(\theta\)，则
\[
\sin\theta=\frac{\left|\overrightarrow{BC_{1}}\cdot\bs{n}\right|}{\left|\overrightarrow{BC_{1}}\right|\left|\bs{n}\right|}
 =\frac{\frac{\sqrt3}{2}}{\sqrt3}=\frac12
\]
因为 \(0\leqslant\theta\leqslant\dfrac\pi2\)，所以 \(\theta=\dfrac\pi6\).
\end{solution}

\begin{problem}
已知 \(\odot O\) 的方程是 \(x^{2} + y^{2} - 2 = 0\)，\(\odot O'\) 的方程是 \(x^{2} + y^{2} - 8x + 10 = 0\)，由动点 \(P\) 向 \(\odot O\) 和 \(\odot O'\) 所引的切线长相等，则动点 \(P\) 的轨迹方程是\fillinblank{}.
\end{problem}

\begin{answer}
\(x=\dfrac32\)
\end{answer}

\begin{solution}
设动点 \(P(x,y)\) 向两圆所引的切线长分别为 \(L_{1}\)，\(L_{2}\). 将两圆方程化为标准形式，得
\[
\odot O:x^{2}+y^{2}=2
\]
\[
\odot O':(x-4)^{2}+y^{2}=6
\]
由圆的切线长公式，
\[
L_{1}^{2}=x^{2}+y^{2}-2
\]
\[
L_{2}^{2}=(x-4)^{2}+y^{2}-6
\]
由 \(L_{1}=L_{2}\)，得到
\[
x^{2}+y^{2}-2=(x-4)^{2}+y^{2}-6=x^{2}+y^{2}-8x+10
\]
所以 \(8x=12\)，动点 \(P\) 的轨迹方程为 \(x=\dfrac32\).
当 \(x=\dfrac32\) 时，\(L_{1}^{2}=L_{2}^{2}=y^{2}+\dfrac14>0\)，故直线上的每一点均在两圆外且能作出两圆的切线，且两切线长相等，所以轨迹确为整条直线 \(x=\dfrac32\).
\end{solution}

\begin{problem}
下面有五个命题：

\begin{circlelist}
\item 函数 \( y = \sin^4 x - \cos^4 x \) 的最小正周期是 \( \pi \)；
\item 终边在 \( y \) 轴上的角的集合是 \( \left\{ \alpha \mid \alpha = \frac{k\pi}{2}, k \in \Z \right\} \)；
\item 在同一坐标系中，函数 \( y = \sin x \) 的图象和函数 \( y = x \) 的图象有 \(3\) 个公共点；
\item 把函数 \( y = 3\sin \left(2x + \frac{\pi}{3}\right) \) 的图象向右平移 \( \frac{\pi}{6} \) 得到 \( y = 3\sin 2x \) 的图象；
\item 角 \(\theta\) 为第一象限角的充要条件是 \(\sin \theta > 0\).
\end{circlelist}

其中真命题的序号是\fillinblank{}.（写出所有）
\end{problem}

\begin{answer}
\(\circled{1}\)\(\circled{4}\)
\end{answer}

\begin{solution}
逐一判断这五个命题. 有
\[
\sin^{4}x-\cos^{4}x=(\sin^{2}x-\cos^{2}x)(\sin^{2}x+\cos^{2}x)
 =\sin^{2}x-\cos^{2}x=-\cos2x
\]
函数 \(-\cos2x\) 的最小正周期为 \(\pi\)，所以命题 \(\circled{1}\) 正确.

终边在 \(y\) 轴上的角应满足
\(\alpha=\frac{(2k+1)\pi}{2}\)，\(k\in\Z\)
而命题 \(\circled{2}\) 还包含终边在 \(x\) 轴上的角，所以命题 \(\circled{2}\) 错误.

令 \(h(x)=x-\sin x\)，则 \(h'(x)=1-\cos x\geqslant0\). 当 \(x>0\) 时，\(h'\) 在 \([0,x]\) 上不恒为零，故 \(h(x)>h(0)=0\)；当 \(x<0\) 时，\(h\) 为奇函数，故 \(h(x)<0\). 因此方程 \(\sin x=x\) 只有 \(x=0\) 一个解，命题 \(\circled{3}\) 错误.

函数 \(y=3\sin\left(2x+\dfrac\pi3\right)\) 的图象向右平移 \(\dfrac\pi6\) 后，得到
\[
y=3\sin\left(2\left(x-\frac\pi6\right)+\frac\pi3\right)=3\sin2x
\]
所以命题 \(\circled{4}\) 正确.

命题 \(\circled{5}\) 不正确，例如 \(\theta=\dfrac{3\pi}{4}\) 时 \(\sin\theta>0\)，但 \(\theta\) 不是第一象限角.

因此真命题的序号为 \(\circled{1}\)\(\circled{4}\).
\end{solution}

\section{解答题}
\begin{problem}
厂家在产品出厂前，需对产品做检验. 厂家对一批产品发给商家的，商家按合同规定抽取一定数量的产品做检验，以决定是否验收这些产品.
\begin{enumerate}
\item 若厂家库房中的每件产品合格的概率为 \(0.3\)，从中任意取出 \(4\) 件进行检验，求至少有 \(1\) 件是合格产品的概率；
\item 若厂家发给商家 \(20\) 件产品，其中有 \(3\) 件不合格，按合同规定该商家从中任取 \(2\) 件，来进行检验，只有 \(2\) 件产品合格时才接收这些产品. 否则拒收，分别求出该商家计算出不合格产品为 \(1\) 件和 \(2\) 件的概率，并求该商家拒收这些产品的概率.
\end{enumerate}

\end{problem}

\begin{solution}
\begin{enumerate}
\item 设抽取的 \(4\) 件产品中合格产品的件数为随机变量 \(X\). 按独立重复抽取的概率模型，每件产品合格的概率为 \(0.3\)，不合格的概率为 \(0.7\)，所以
\[
P(X\geqslant1)=1-P(X=0)=1-(1-0.3)^4=1-0.7^4=0.7599
\]
\item 这 \(20\) 件产品中有 \(17\) 件合格，任取 \(2\) 件的基本事件数为 \(\mathrm C_{20}^{2}\). 设抽到的不合格产品数为随机变量 \(Y\)，则
\[
P(Y=1)=\frac{\mathrm C_{3}^{1}\mathrm C_{17}^{1}}{\mathrm C_{20}^{2}}=\frac{51}{190}
\]
\[
P(Y=2)=\frac{\mathrm C_{3}^{2}}{\mathrm C_{20}^{2}}=\frac{3}{190}
\]
只有抽到 \(2\) 件合格产品时才接收，所以拒收的事件为 \(Y\geqslant1\). 因而
\[
P(\text{拒收})=P(Y=1)+P(Y=2)=\frac{51+3}{190}=\frac{27}{95}
\]
\end{enumerate}
\end{solution}

\begin{problem}
已知 \(\cos \alpha = \frac{1}{7}\)，\(\cos (\alpha - \beta) = \frac{13}{14}\)，且 \(0 < \beta < \alpha < \frac{\pi}{2}\).
\begin{enumerate}
\item 求 \( \tan 2\alpha \) 的值；
\item 求 \( \beta \).
\end{enumerate}

\end{problem}

\begin{solution}
\begin{enumerate}
\item 因为 \(0<\alpha<\dfrac\pi2\)，所以 \(\sin\alpha=\sqrt{1-\cos^2\alpha}=\dfrac{4\sqrt3}{7}\)，从而 \(\tan\alpha=4\sqrt3\). 因此
\[
\tan2\alpha=\frac{2\tan\alpha}{1-\tan^2\alpha}
 =\frac{8\sqrt3}{1-48}=-\frac{8\sqrt3}{47}
\]
\item 由 \(0<\alpha-\beta<\dfrac\pi2\)，得
\[
\sin(\alpha-\beta)=\sqrt{1-\cos^2(\alpha-\beta)}=\frac{3\sqrt3}{14}
\]
利用差角公式，
\[
\cos\beta=\cos\bigl(\alpha-(\alpha-\beta)\bigr)
 =\cos\alpha\cos(\alpha-\beta)+\sin\alpha\sin(\alpha-\beta)
 =\frac{13}{98}+\frac{36}{98}=\frac12
\]
又 \(0<\beta<\dfrac\pi2\)，所以 \(\beta=\dfrac\pi3\).
\end{enumerate}
\end{solution}

\begin{problem}
如图，平面 \(PCBM\perp\) 平面 \(ABC\)，\(PCBM\) 是直角梯形，\(\angle PCB = 90^{\circ}\)，\(PM \parallel BC\)，直线 \(AM\) 与直线 \(PC\) 所成的角为 \(60^{\circ}\)，又 \(AC = 1\)，\(BC = 2PM = 2\)，\(\angle ACB = 90^{\circ}\).
\begin{enumerate}
\item 求证：\( AC \perp BM \)；
\item 求二面角 \(M - AB - C\) 的大小；
\item 求多面体 \(PMABC\) 的体积.
\end{enumerate}

\bitmapfigure[max width=6.0cm]{img/2007/sichuan_liberal/q19_fig1.png}
\end{problem}

\begin{solution}
建立空间直角坐标系，使平面 \(ABC\) 为 \(xOy\) 平面，取 \(C=(0,0,0)\)，\(B=(2,0,0)\)，\(A=(0,1,0)\). 由平面 \(PCBM\perp\) 平面 \(ABC\)，且两平面的交线为 \(BC\)，可令平面 \(PCBM\) 为 \(xOz\) 平面. 由 \(\angle PCB=90^{\circ}\)，设 \(P=(0,0,h)\)，其中 \(h>0\). 由于 \(PM\parallel BC\)，且 \(BC=2PM=2\)，得 \(PM=1\)，从而 \(M=(1,0,h)\).

\begin{enumerate}
\item 取直线 \(PC\) 的方向向量 \(\bs{u}=(0,0,h)\)，直线 \(AM\) 的方向向量 \(\bs{v}=(1,-1,h)\). 由两直线所成的角为 \(60^{\circ}\)，有
\[
\cos60^{\circ}=\frac{|\bs{u}\cdot\bs{v}|}{|\bs{u}|\,|\bs{v}|}=\frac{h^{2}}{h\sqrt{h^{2}+2}}=\frac{h}{\sqrt{h^{2}+2}}
\]
因此 \(4h^{2}=h^{2}+2\)，即 \(h^{2}=\dfrac23\).

又 \(\overrightarrow{AC}=(0,-1,0)\)，\(\overrightarrow{BM}=(-1,0,h)\)，所以
\[
\overrightarrow{AC}\cdot\overrightarrow{BM}=0
\]
故 \(AC\perp BM\).
\item \(\overrightarrow{AB}=(2,-1,0)\)，\(\overrightarrow{AM}=(1,-1,h)\). 取平面 \(MAB\) 的法向量
\[
\bs{n}_{1}=\overrightarrow{AB}\times\overrightarrow{AM}=(-h,-2h,-1)
\]
取平面 \(ABC\) 的法向量 \(\bs{n}_{2}=(0,0,1)\). 设二面角 \(M-AB-C\) 的大小为 \(\theta\)，则
\[
\cos\theta=\frac{|\bs{n}_{1}\cdot\bs{n}_{2}|}{|\bs{n}_{1}|\,|\bs{n}_{2}|}=\frac{1}{\sqrt{5h^{2}+1}}=\frac{1}{\sqrt{\frac{13}{3}}}=\frac{\sqrt{39}}{13}
\]
所以 \(\theta=\arccos\dfrac{\sqrt{39}}{13}\).
\item 四边形 \(PCBM\) 是以 \(BC\)，\(PM\) 为两底的直角梯形，且高为 \(PC=h\)，所以
\[
S_{PCBM}=\frac{BC+PM}{2}\cdot PC=\frac{2+1}{2}h=\frac32h
\]
由平面 \(PCBM\perp\) 平面 \(ABC\)，且 \(AC\perp BC\)，可知 \(AC\perp\) 平面 \(PCBM\). 因此以四边形 \(PCBM\) 为底、点 \(A\) 为顶点的多面体 \(PMABC\) 的高为 \(AC=1\)，从而
\[
V_{PMABC}=\frac13S_{PCBM}\cdot AC=\frac13\cdot\frac32h\cdot1=\frac h2=\frac12\sqrt{\frac23}=\frac{\sqrt6}{6}
\]
\end{enumerate}
\end{solution}

\begin{problem}
设函数 \( f(x) = ax^3 + bx + c \)（\(a \neq 0\)）为奇函数，其图象在点 \((1, f(1))\) 处的切线与直线 \( x - 6y - 7 = 0 \) 垂直，导函数 \( f'(x) \) 的最小值为 \(-12\).
\begin{enumerate}
\item 求 \(a\)，\(b\)，\(c\) 的值；
\item 求函数 \( f(x) \) 的单调递增区间，并求函数 \( f(x) \) 在 \( [-1,3] \) 上的最大值和最小值.
\end{enumerate}

\end{problem}

\begin{solution}
\begin{enumerate}
\item 由 \(f(x)=ax^{3}+bx+c\) 为奇函数，得
\[
f(-x)=-ax^{3}-bx+c=-f(x)=-ax^{3}-bx-c
\]
所以 \(c=0\). 又直线 \(x-6y-7=0\) 的斜率为 \(\dfrac16\)，故点 \((1,f(1))\) 处切线的斜率为 \(-6\). 因此
\[
f'(1)=3a+b=-6
\]
而 \(f'(x)=3ax^{2}+b\). 若 \(a<0\)，则 \(f'(x)\) 没有最小值；所以 \(a>0\)，且 \(f'(x)\) 的最小值为 \(f'(0)=b=-12\). 代入 \(3a+b=-6\)，得 \(a=2\). 因此
\[
(a,b,c)=(2,-12,0)
\]
且 \(f(x)=2x^{3}-12x\).
\item \(f'(x)=6x^{2}-12=6(x^{2}-2)\). 所以当 \(x<-\sqrt2\) 或 \(x>\sqrt2\) 时，\(f'(x)>0\)，函数单调递增；当 \(-\sqrt2<x<\sqrt2\) 时，\(f'(x)<0\). 因此函数的单调递增区间为 \((-\infty,-\sqrt2)\) 和 \((\sqrt2,+\infty)\). 在区间 \([-1,3]\) 内，只有临界点 \(x=\sqrt2\). 计算端点和临界点处的函数值：
计算得 \(f(-1)=10\)，\(f(\sqrt2)=-8\sqrt2\)，\(f(3)=18\). 因为 \(-8\sqrt2<10<18\)，所以函数在 \([-1,3]\) 上的最大值为 \(18\)，最小值为 \(-8\sqrt2\).
\end{enumerate}
\end{solution}

\begin{problem}
已知 \(F_{1}\)，\(F_{2}\) 分别是椭圆 \(\frac{x^{2}}{4} + y^{2} = 1\) 的左，右焦点.
\begin{enumerate}
\item 若 \(P\) 是第一象限内该椭圆上的一点，且 \(\overrightarrow{PF_{1}} \cdot \overrightarrow{PF_{2}} = -\frac{5}{4}\)，求点 \(P\) 的坐标；
\item 设过定点 \(M(0,2)\) 的直线 \(l\) 与椭圆交于不同的两点 \(A\)，\(B\)，且 \(\angle AOB\) 为锐角（其中 \(O\) 为坐标原点），求直线 \(l\) 的斜率 \(k\) 的取值范围.
\end{enumerate}

\end{problem}

\begin{solution}
椭圆的长半轴为 \(a=2\)，短半轴为 \(b=1\)，所以焦半距
\[
c=\sqrt{a^{2}-b^{2}}=\sqrt3
\]
从而 \(F_{1}=(-\sqrt3,0)\)，\(F_{2}=(\sqrt3,0)\).

\begin{enumerate}
\item 设第一象限内的点 \(P\) 为 \((x,y)\)，则 \(x>0\)，\(y>0\)，且
\[
\frac{x^{2}}4+y^{2}=1
\]
又
\(\overrightarrow{PF_{1}}=(-\sqrt3-x,-y)\)，\(\overrightarrow{PF_{2}}=(\sqrt3-x,-y)\)，
所以
\[
\overrightarrow{PF_{1}}\cdot\overrightarrow{PF_{2}}
=(-\sqrt3-x)(\sqrt3-x)+y^{2}=x^{2}+y^{2}-3
\]
由题意，\(x^{2}+y^{2}-3=-\dfrac54\)，即 \(x^{2}+y^{2}=\dfrac74\). 与椭圆方程相减，得
\[
\frac34x^{2}=\frac34
\]
所以 \(x^{2}=1\). 因为 \(x>0\)，故 \(x=1\)，再由椭圆方程得 \(y=\dfrac{\sqrt3}{2}\). 因此 \(P\left(1,\dfrac{\sqrt3}{2}\right)\).
\item 先排除竖直直线 \(x=0\)：它与椭圆交于 \((0,1)\)，\((0,-1)\)，两交点的位置向量方向相反，故 \(\angle AOB=\pi\)，不是锐角. 所以符合条件的直线可设斜率为 \(k\)，方程为 \(y=kx+2\). 代入椭圆方程，得
\[
\left(k^{2}+\frac14\right)x^{2}+4kx+3=0
\]
两交点不同的充要条件是判别式大于零，即
\[
\Delta=(4k)^{2}-12\left(k^{2}+\frac14\right)=4k^{2}-3>0
\]
所以 \(\lvert k\rvert>\dfrac{\sqrt3}{2}\).

设两交点的横坐标为 \(x_{1}\)，\(x_{2}\). 由韦达定理，
\[
x_{1}+x_{2}=-\frac{16k}{4k^{2}+1}
\]
\[
x_{1}x_{2}=\frac{12}{4k^{2}+1}
\]
两交点的位置向量的内积为
\[
\overrightarrow{OA}\cdot\overrightarrow{OB}
=x_{1}x_{2}+(kx_{1}+2)(kx_{2}+2)
=\frac{4(4-k^{2})}{4k^{2}+1}
\]
由于 \(\angle AOB\) 为锐角，当且仅当该内积大于零，即 \(\lvert k\rvert<2\). 综合两交点不同的条件，得到
\[
k\in\left(-2,-\frac{\sqrt3}{2}\right)\cup\left(\frac{\sqrt3}{2},2\right)
\]
\end{enumerate}
\end{solution}

\begin{problem}
已知函数 \( f(x) = x^2 - 4 \)，设曲线 \( y = f(x) \) 在点 \( (x_n, f(x_n)) \) 处的切线与 \( x \) 轴的交点为 \( (x_{n+1}, 0) \)（\(n \in \mathbf{N}^*\)），其中 \( x_1 \) 为正实数.
\begin{enumerate}
\item 用 \(x_{n}\) 表示 \(x_{n + 1}\)；
\item 若 \(x_{1} = 4\)，记 \(a_{n} = \lg \frac{x_{n} + 2}{x_{n} - 2}\)，证明数列 \(\{a_{n}\}\) 成等比数列，并求数列 \(\{x_{n}\}\) 的通项公式；
\item 若 \(x_{1} = 4\)，\(b_{n} = x_{n} - 2\)，\(T_{n}\) 是数列 \(\{b_{n}\}\) 的前 \(n\) 项和，证明 \(T_{n} < 3\).
\end{enumerate}

\end{problem}

\begin{solution}
\begin{enumerate}
\item 曲线在点 \((x_{n},x_{n}^{2}-4)\) 处的切线斜率为 \(2x_{n}\)，故切线方程为
\[
y-(x_{n}^{2}-4)=2x_{n}(x-x_{n})
\]
令 \(y=0\)，并令横坐标为 \(x_{n+1}\)，得
\(4-x_{n}^{2}=2x_{n}(x_{n+1}-x_{n})\)，\(2x_{n}x_{n+1}=x_{n}^{2}+4\).
因为 \(x_{1}>0\)，若 \(x_{n}>0\)，则由上式可得
\[
x_{n+1}=\frac{x_{n}^{2}+4}{2x_{n}}>0
\]
由归纳法，所有 \(x_{n}>0\)，从而
\[
x_{n+1}=\frac{x_{n}}2+\frac{2}{x_{n}}
\]
\item 当 \(x_{1}=4\) 时，先有 \(x_{1}>2\). 若 \(x_{n}>2\)，则
\[
x_{n+1}-2=\frac{(x_{n}-2)^{2}}{2x_{n}}>0
\]
故归纳可知对任意正整数 \(n\)，都有 \(x_{n}>2\)，从而下式中的对数有意义. 又
\(x_{n+1}+2=\frac{(x_{n}+2)^{2}}{2x_{n}}\)，\(x_{n+1}-2=\frac{(x_{n}-2)^{2}}{2x_{n}}\)
因此
\[
\frac{x_{n+1}+2}{x_{n+1}-2}=\left(\frac{x_{n}+2}{x_{n}-2}\right)^{2}
\]
记 \(a_{n}=\lg\dfrac{x_{n}+2}{x_{n}-2}\)，则 \(a_{n+1}=2a_{n}\)，且
\[
a_{1}=\lg\frac{4+2}{4-2}=\lg3
\]
所以 \(\{a_{n}\}\) 是首项为 \(\lg3\)、公比为 \(2\) 的等比数列，且
\[
a_{n}=2^{n-1}\lg3
\]
令 \(r_{n}=10^{a_{n}}=3^{2^{n-1}}\)，则 \(\dfrac{x_{n}+2}{x_{n}-2}=r_{n}\)，从而
\[
x_{n}=\frac{2(r_{n}+1)}{r_{n}-1}
=\frac{2\left(3^{2^{n-1}}+1\right)}{3^{2^{n-1}}-1}
\]
\item 由第（1）问，令 \(b_{n}=x_{n}-2\)，则
\[
b_{n+1}=x_{n+1}-2=\frac{b_{n}^{2}}{2(b_{n}+2)}
\]
其中 \(b_{1}=2\)，所以 \(b_{2}=\dfrac12\). 若 \(0<b_{n}\leqslant\dfrac12\)，则
\[
0<\frac{b_{n+1}}{b_{n}}=\frac{b_{n}}{2(b_{n}+2)}\leqslant\frac1{10}<\frac14
\]
从而 \(0<b_{n+1}<\dfrac12\). 由归纳法，对 \(n\geqslant2\) 均有 \(0<b_{n}\leqslant\dfrac12\)，并且
\[
b_{n}\leqslant\frac12\left(\frac14\right)^{n-2}
\]
当 \(n=1\) 时，\(T_{1}=b_{1}=2<3\)；当 \(n\geqslant2\) 时，
\[
T_{n}=b_{1}+\sum_{j=2}^{n}b_{j}
\leqslant2+\frac12\sum_{j=0}^{n-2}\left(\frac14\right)^{j}
<2+\frac12\cdot\frac{1}{1-\frac14}=\frac83<3
\]
故对任意正整数 \(n\)，均有 \(T_{n}<3\).
\end{enumerate}
\end{solution}
